\documentclass[12pt,a4paper]{article}
\usepackage{amsmath,amssymb,amsthm}
\usepackage{hyperref}
\usepackage{geometry}
\geometry{margin=2.5cm}
\usepackage{enumitem}
\usepackage{booktabs}

\newtheorem{theorem}{Theorem}[section]
\newtheorem{lemma}[theorem]{Lemma}
\newtheorem{corollary}[theorem]{Corollary}
\newtheorem{proposition}[theorem]{Proposition}
\theoremstyle{definition}
\newtheorem{definition}[theorem]{Definition}
\theoremstyle{remark}
\newtheorem{remark}[theorem]{Remark}

\title{The Core-Collapse Supernova Mechanism\\from Recognition Science}

\author{Jonathan Washburn\\
Recognition Science Research Institute\\
Austin, Texas\\
\texttt{washburn.jonathan@gmail.com}}

\date{March 2026}

\begin{document}
\maketitle

\begin{abstract}
We analyze the core-collapse supernova mechanism within Recognition
Science (RS).  The collapse is triggered when the iron core of a
massive star ($M \gtrsim 8\,M_\odot$) exceeds the Chandrasekhar mass
$M_{\mathrm{Ch}} \approx 1.44\,M_\odot$ (derived from RS in a
companion paper).  The bounce at nuclear density
($\rho_{\mathrm{nuc}} \approx 2.8 \times 10^{14}\;\mathrm{g/cm^3}$)
is a ledger phase transition: the recognition events become maximally
compressed ($\ell \to \ell_0$), halting the collapse.  The explosion
is driven by neutrino energy deposition, with the neutrino
interaction cross-sections determined by the RS-derived electroweak
parameters.  The total energy release $E_{\mathrm{tot}} \sim
3 \times 10^{53}\;\mathrm{erg}$ is the gravitational binding energy
of the neutron star remnant, of which ${\sim}\,99\%$ is carried by
neutrinos.  The $\varphi$-ladder structure predicts the neutron star
mass distribution and a maximum mass for neutron stars.
\end{abstract}

%% ============================================================
\section{Introduction}
\label{sec:intro}
%% ============================================================

Core-collapse supernovae are among the most energetic events in the
universe~\cite{Janka2012}.  Despite decades of numerical simulation,
the explosion mechanism remains not fully understood: 1D simulations
fail to explode, and 3D simulations succeed only with careful
treatment of neutrino transport and multi-dimensional
instabilities~\cite{Burrows2021}.  Recognition Science~\cite{RS_Axioms}
provides a structural framework for understanding each phase of the
collapse and explosion.

%% ============================================================
\section{Collapse Trigger}
\label{sec:trigger}
%% ============================================================

\begin{theorem}[Chandrasekhar Mass]
\label{thm:chandra}
The maximum mass of a degenerate electron core (derived from the
RS $\varphi$-ladder in a companion paper~\cite{RS_Chandrasekhar}):
\begin{equation}
  M_{\mathrm{Ch}} = 5.83\,\mu_e^{-2}\,M_\odot
  \approx 1.44\,M_\odot,
\end{equation}
where $\mu_e \approx 2$ for compositions with $Y_e \approx 0.5$.
For the pre-collapse iron core ($Y_e \approx 0.42$,
$\mu_e \approx 2.15$), electron captures reduce the
effective Chandrasekhar mass to
$M_{\mathrm{Ch}}^{\mathrm{eff}} \approx 1.2$--$1.3\,M_\odot$.
\end{theorem}

\begin{theorem}[Collapse Initiation]
When $M_{\mathrm{core}} > M_{\mathrm{Ch}}$, electron degeneracy
pressure can no longer support the core.  Electron capture
($p + e^- \to n + \nu_e$) further reduces $Y_e$ and the pressure
support, triggering dynamical collapse on a free-fall timescale:
\begin{equation}
  t_{\mathrm{ff}} = \sqrt{\frac{3\pi}{32\,G\,\rho}} \approx
  0.1\;\mathrm{s}.
\end{equation}
\end{theorem}

%% ============================================================
\section{Bounce and Shock Formation}
\label{sec:bounce}
%% ============================================================

\begin{theorem}[Nuclear Density Bounce]
\label{thm:bounce}
The core collapse halts at nuclear density
$\rho_{\mathrm{nuc}} \approx 2.8 \times
10^{14}\;\mathrm{g/cm^3}$.
\end{theorem}

\begin{proof}
At nuclear density, the recognition events in the core are maximally
compressed: the inter-voxel spacing approaches $\ell_0$, the minimum
resolution of the ledger.  The nuclear equation of state (EOS)
stiffens dramatically due to the strong nuclear force (repulsive at
short range), creating an effectively incompressible state.  The
infalling material bounces off this hard core, launching a shock
wave.
\end{proof}

\begin{definition}[Proto-Neutron Star]
The hot, lepton-rich remnant formed at bounce, with:
\begin{align}
  M_{\mathrm{PNS}} &\approx 0.6\,M_\odot
  \quad (\text{initial}), \\
  R_{\mathrm{PNS}} &\approx 50\;\mathrm{km}, \\
  T_{\mathrm{PNS}} &\approx 30\;\mathrm{MeV}.
\end{align}
\end{definition}

%% ============================================================
\section{Neutrino-Driven Explosion}
\label{sec:neutrinos}
%% ============================================================

\begin{theorem}[Shock Stalling and Revival]
The bounce shock stalls at $R \approx 150$--$200\;\mathrm{km}$ due
to energy losses from photodissociation of iron nuclei and neutrino
emission.  The shock is revived by neutrino energy deposition in the
``gain region'' between the neutrinosphere and the stalled shock.
\end{theorem}

\begin{theorem}[Energy Budget]
\label{thm:energy}
\begin{align}
  E_{\mathrm{tot}} &\approx \frac{3\,G\,M_{\mathrm{NS}}^2}
  {5\,R_{\mathrm{NS}}} \approx 3 \times
  10^{53}\;\mathrm{erg}, \\
  E_\nu &\approx 0.99 \times E_{\mathrm{tot}}, \\
  E_{\mathrm{kin}} &\approx 10^{51}\;\mathrm{erg}
  \approx 1\;\text{Bethe (B)}, \\
  E_{\mathrm{photon}} &\approx 10^{49}\;\mathrm{erg}.
\end{align}
\end{theorem}

\begin{proof}
The gravitational binding energy of a uniform-density neutron star
with $M = 1.4\,M_\odot$ and $R = 10\;\mathrm{km}$ is
$E_{\mathrm{grav}} = 3GM^2/(5R) \approx 3 \times
10^{53}\;\mathrm{erg}$.  In RS, this binding energy is computed with
$G = \varphi^5/\pi$ in natural units; the numerical result is
identical.  The neutrino luminosity dominates because neutrinos
decouple at the neutrinosphere (optical depth $\tau_\nu \sim 1$ at
$R \sim 10$--$50\;\mathrm{km}$), carrying away the thermal energy.
\end{proof}

%% ============================================================
\section{Neutrino Cross-Sections}
\label{sec:cross-sections}
%% ============================================================

\begin{theorem}[Neutrino Absorption]
The charged-current neutrino absorption cross-section:
\begin{equation}
  \sigma(\nu_e + n \to p + e^-) \approx \frac{G_F^2\,E_\nu^2}
  {\pi} \approx 9 \times 10^{-44}\;\mathrm{cm^2}
  \left(\frac{E_\nu}{10\;\mathrm{MeV}}\right)^2,
\end{equation}
where $G_F$ is the Fermi constant (derived from the RS electroweak
scale $v = m_e \varphi^{27}$).
\end{theorem}

\begin{remark}
The neutrino heating rate in the gain region:
\begin{equation}
  \dot{Q}_\nu \sim \frac{L_\nu\,\sigma}{4\pi R^2\,m_N}
  \propto L_\nu\,E_\nu^2\,R^{-2},
\end{equation}
must exceed the cooling rate for successful explosion.  This
condition is sensitive to the neutrino luminosity $L_\nu$ and mean
energy $E_\nu$, both determined by the RS-derived electroweak
parameters and nuclear equation of state.
\end{remark}

%% ============================================================
\section{Neutron Star Mass Distribution}
\label{sec:ns-mass}
%% ============================================================

\begin{theorem}[Neutron Star Mass Range]
RS predicts a neutron star mass distribution peaked at
$M_{\mathrm{NS}} \approx 1.4\,M_\odot$ with a maximum mass
$M_{\mathrm{max}} \approx 2.0$--$2.3\,M_\odot$.
\end{theorem}

\begin{remark}
The observed distribution~\cite{Ozel2016} peaks at $1.33\,M_\odot$
(for double neutron star systems) and extends to
${\sim}\,2.1\,M_\odot$ (PSR~J0740+6620).  The RS prediction is
consistent.
\end{remark}

%% ============================================================
\section{Predictions}
\label{sec:predictions}
%% ============================================================

\begin{enumerate}
  \item \textbf{SN~1987A neutrinos confirmed the energy budget}:
  The ${\sim}\,20$ neutrino events detected by Kamiokande-II and
  IMB~\cite{Hirata1987} are consistent with
  $E_\nu \sim 3 \times 10^{53}\;\mathrm{erg}$ and mean energy
  $\langle E_\nu \rangle \sim 12\;\mathrm{MeV}$.

  \item \textbf{Neutrino mass ordering}: The $\varphi$-ladder
  neutrino mass hierarchy predicts detectable spectral differences
  between $\nu_e$, $\bar{\nu}_e$, and $\nu_x$ in the next Galactic
  supernova.

  \item \textbf{Neutron star maximum mass}: $M_{\mathrm{max}}
  \lesssim 2.3\,M_\odot$.  Any neutron star exceeding this would
  challenge the RS nuclear equation of state.
\end{enumerate}

%% ============================================================
\section{Conclusion}
\label{sec:conclusion}
%% ============================================================

The core-collapse supernova mechanism in Recognition Science
proceeds through RS-derived physics at every stage: the
Chandrasekhar mass triggers collapse, the ledger voxel structure
provides the bounce at nuclear density, and the RS-derived
electroweak parameters determine the neutrino cross-sections
that drive the explosion.  The total energy budget of
${\sim}\,3 \times 10^{53}\;\mathrm{erg}$ is computed with zero
free parameters.

%% ============================================================
\begin{thebibliography}{99}

\bibitem{RS_Axioms}
J.~Washburn, ``Recognition Science: Foundations,''
RS Axioms Document (2026).

\bibitem{RS_Chandrasekhar}
J.~Washburn, ``The Chandrasekhar Mass from Recognition Science,''
RS papers (2026).

\bibitem{Janka2012}
H.-T.~Janka, ``Explosion Mechanisms of Core-Collapse Supernovae,''
Ann.\ Rev.\ Nucl.\ Part.\ Sci.\ \textbf{62}, 407 (2012).

\bibitem{Burrows2021}
A.~Burrows and D.~Vartanyan, ``Core-Collapse Supernova Explosion
Theory,'' Nature \textbf{589}, 29 (2021).

\bibitem{Ozel2016}
F.~\"Ozel and P.~Freire, ``Masses, Radii, and the Equation of State
of Neutron Stars,'' Ann.\ Rev.\ Astron.\ Astrophys.\ \textbf{54},
401 (2016).

\bibitem{Hirata1987}
K.~Hirata \textit{et al.}, ``Observation of a Neutrino Burst from
the Supernova SN~1987A,'' Phys.\ Rev.\ Lett.\ \textbf{58}, 1490
(1987).

\end{thebibliography}

\end{document}
